\documentclass[10pt]{beamer}

% Tema Limpo e Profissional
\usetheme{Madrid}
\usecolortheme{whale}

% Pacotes Essenciais
\usepackage[utf8]{inputenc}
\usepackage[portuguese]{babel}
\usepackage{graphicx}
\usepackage{booktabs}
\usepackage{amsmath}
\usepackage{amssymb}
\usepackage{listings}
\usepackage{xcolor}
\usepackage{tikz}
\usetikzlibrary{arrows.meta, positioning, shapes.geometric}

% Configuracoes de Codigo (Python)
\lstset{
    language=Python,
    basicstyle=\ttfamily\scriptsize,
    keywordstyle=\color{blue}\bfseries,
    stringstyle=\color{orange},
    commentstyle=\color{gray!80!black}\itshape,
    breaklines=true,
    showstringspaces=false,
    backgroundcolor=\color{gray!5},
    frame=single,
    rulecolor=\color{gray!30},
    numbers=left,
    numberstyle=\tiny\color{gray},
    stepnumber=1,
    numbersep=5pt
}

% Metadados da Apresentacao
\title[Pygame - Introdução 2D]{Introdução ao Desenvolvimento de Jogos com Pygame}
\subtitle{Arquitetura do Game Loop, Coordenadas 2D e Interatividade}
\author[Prof. Reginaldo]{Prof. Reginaldo Fernandes\texorpdfstring{\\ \small{Programação de Computadores / ADS}}{}}
\institute[IFCE]{Instituto Federal do Ceará -- Campus Tauá \\ Tecnologia em Análise e Desenvolvimento de Sistemas}
\date{\today}

\begin{document}

% Slide de Titulo
\begin{frame}
    \titlepage
\end{frame}

% Sumario
\begin{frame}{Sumário da Aula (Duração: 90 min)}
    \tableofcontents
\end{frame}

% Blocos Modulares de Slides
\section{Introdução e Filosofia do Pygame}

\begin{frame}{O que é o Pygame?}
    \begin{itemize}
        \item O \textbf{Pygame} é um conjunto de módulos em Python projetados para a criação de jogos bidimensionais (2D) e aplicações multimídia.
        \item Ele é construído sobre a biblioteca em C \textbf{SDL} (\textit{Simple DirectMedia Layer}), combinando a agilidade do Python com a performance do hardware gráfico.
    \end{itemize}

    \vspace{0.3cm}
    \begin{columns}
        \begin{column}{0.5\textwidth}
            \textbf{O que o Pygame gerencia?}
            \begin{itemize}
                \item Criação de janelas gráficas.
                \item Renderização de formas e imagens.
                \item Captura de teclado, mouse e controles.
                \item Reprodução de sons e efeitos de áudio.
                \item Controle fino do relógio (FPS).
            \end{itemize}
        \end{column}
        \begin{column}{0.5\textwidth}
            \begin{block}{Instalação Simples}
                No terminal do laboratório:
                \begin{center}
                    \texttt{pip install pygame-ce}
                \end{center}
                \small{\textit{pygame-ce} (Community Edition) é a versão moderna mantida pela comunidade.}
            \end{block}
        \end{column}
    \end{columns}
\end{frame}

\begin{frame}{Programas Tradicionais vs.\ Jogos em Tempo Real}
    \begin{columns}
        \begin{column}{0.48\textwidth}
            \begin{block}{Aplicações Tradicionais (Web/CLI)}
                \begin{itemize}
                    \item Modelo \textbf{orientado a eventos passivo}.
                    \item O programa fica \textbf{em repouso} aguardando o usuário digitar ou clicar.
                    \item Se ninguém interagir, nada acontece na tela.
                \end{itemize}
            \end{block}
        \end{column}
        \begin{column}{0.48\textwidth}
            \begin{block}{Jogos em Tempo Real}
                \begin{itemize}
                    \item Modelo de \textbf{Laço Contínuo} (\textit{Game Loop}).
                    \item O mundo virtual continua acontecendo \textbf{60 vezes por segundo}.
                    \item Inimigos andam, física se aplica e a tela redesenha, mesmo sem comandos.
                \end{itemize}
            \end{block}
        \end{column}
    \end{columns}

    \vspace{0.4cm}
    \begin{center}
        \textbf{Conclusão Didática:} Um jogo é essencialmente uma animação interativa calculada em tempo real dentro de um laço de repetição infinito!
    \end{center}
\end{frame}

\section{A Anatomia do Game Loop}

\begin{frame}{O Ciclo de Vida de um Jogo}
    Todo jogo de computador, do clássico Pong aos títulos AAA modernos, é executado dentro de quatro fases essenciais a cada quadro (\textit{frame}):

    \vspace{0.3cm}
    \begin{center}
    \begin{tikzpicture}[
        node distance=1.0cm and 1.5cm,
        box/.style={rectangle, draw=blue!70, fill=blue!10, rounded corners, text centered, minimum width=2.8cm, minimum height=0.7cm, font=\footnotesize\bfseries},
        arrow/.style={-{Stealth[scale=1.0]}, thick, color=blue!80}
    ]
        \node[box] (eventos) {1. Entrada (Eventos)};
        \node[box, right=of eventos] (logica) {2. Lógica (Update)};
        \node[box, below=of logica] (desenho) {3. Desenho (Render)};
        \node[box, below=of eventos] (clock) {4. Relógio (Clock Tick)};

        \draw[arrow] (eventos) -- (logica);
        \draw[arrow] (logica) -- (desenho);
        \draw[arrow] (desenho) -- (clock);
        \draw[arrow] (clock) -- (eventos);
    \end{tikzpicture}
    \end{center}

    \begin{itemize}
        \item \textbf{1. Entrada:} Lê cliques do mouse, teclas pressionadas e fechamento da janela.
        \item \textbf{2. Lógica:} Atualiza posições, calcula colisões, pontuação e regras.
        \item \textbf{3. Desenho:} Limpa o buffer de vídeo e redesenha todos os elementos.
        \item \textbf{4. Relógio:} Aguarda o tempo necessário para travar a taxa de quadros (ex: 60 FPS).
    \end{itemize}
\end{frame}

\begin{frame}[fragile]{O Código Mínimo Essencial (Esqueleto)}
    \begin{lstlisting}[language=Python]
import pygame, sys

pygame.init()
tela = pygame.display.set_mode((800, 600))
pygame.display.set_caption("Meu Primeiro Jogo")
relogio = pygame.time.Clock()

executando = True
while executando:
    # 1. Eventos
    for evento in pygame.event.get():
        if evento.type == pygame.QUIT:
            executando = False

    # 2. Logica (vazio no momento)

    # 3. Desenho
    tela.fill((30, 30, 45))  # Limpa fundo com cinza escuro
    pygame.display.flip()    # Exibe no monitor

    # 4. Controle de FPS
    relogio.tick(60)

pygame.quit()
sys.exit()
    \end{lstlisting}
\end{frame}

\begin{frame}{Por que o programa trava sem \texttt{pygame.event.get()}?}
    \begin{block}{O Sistema Operacional e a Fila de Mensagens}
        \begin{itemize}
            \item Os sistemas operacionais (Linux, Windows, macOS) enviam mensagens contínuas para as janelas abertas (foco, teclado, redimensionamento, botão de fechar).
            \item Se o jogo não consumir essa fila através de \texttt{pygame.event.get()}, o sistema assume que o processo entrou em congelamento (\textit{loop infinito travado}) e exibe a mensagem de ``Não Respondendo''.
        \end{itemize}
    \end{block}

    \vspace{0.3cm}
    \begin{columns}
        \begin{column}{0.5\textwidth}
            \textbf{O que faz \texttt{pygame.QUIT}?}
            \begin{itemize}
                \item Representa o clique no botão vermelho de fechar janela `X'.
                \item Permite sair suavemente do laço \texttt{while}.
            \end{itemize}
        \end{column}
        \begin{column}{0.5\textwidth}
            \textbf{O que faz \texttt{display.flip()}?}
            \begin{itemize}
                \item Implementa o \textbf{Double Buffering}: tudo é desenhado na memória oculta e depois trocado de uma só vez para a tela visível, evitando cintilação (\textit{flickering}).
            \end{itemize}
        \end{column}
    \end{columns}
\end{frame}

\section{Coordenadas 2D e Cores RGB}

\begin{frame}{O Plano Cartesiano na Computação Gráfica}
    \textbf{Atenção:} As coordenadas de tela na computação gráfica são diferentes do plano cartesiano que aprendemos na matemática tradicional!

    \vspace{0.2cm}
    \begin{columns}
        \begin{column}{0.5\textwidth}
            \begin{center}
            \begin{tikzpicture}[scale=0.65]
                % Moldura da tela
                \draw[thick, fill=gray!10] (0,0) rectangle (5, -4);
                
                % Origem
                \fill[red] (0,0) circle (3pt) node[above right, font=\tiny\bfseries] {(0, 0) Origem};
                
                % Eixo X
                \draw[-{Stealth[scale=1.2]}, blue, very thick] (0,0) -- (5.2, 0) node[right, font=\tiny\bfseries] {+X};
                
                % Eixo Y
                \draw[-{Stealth[scale=1.2]}, red, very thick] (0,0) -- (0, -4.2) node[below, font=\tiny\bfseries] {+Y (Para baixo!)};
                
                % Ponto de exemplo
                \fill[blue!80] (3, -2) circle (3pt) node[right, font=\tiny] {(x=300, y=200)};
                \draw[dashed, gray] (3,0) -- (3,-2) -- (0,-2);
                
                % Borda inferior direita
                \node[above left, font=\tiny] at (5, -4) {(800, 600)};
            \end{tikzpicture}
            \end{center}
        \end{column}
        \begin{column}{0.5\textwidth}
            \begin{block}{Regras Fundamentais}
                \begin{itemize}
                    \item O ponto \textbf{(0, 0)} fica no canto \textbf{superior esquerdo} (\textit{Top-Left}).
                    \item O eixo \textbf{X} cresce para a \textbf{DIREITA}.
                    \item O eixo \textbf{Y} cresce para \textbf{BAIXO}!
                    \item Logo, para descer um objeto na tela, \textbf{somamos} em $Y$; para subir, \textbf{subtraímos} de $Y$.
                \end{itemize}
            \end{block}
        \end{column}
    \end{columns}
\end{frame}

\begin{frame}{O Modelo de Cores Aditivas: RGB}
    As cores no Pygame são representadas por tuplas de 3 números inteiros variando entre $0$ e $255$: \textbf{(Red, Green, Blue)}.

    \vspace{0.3cm}
    \begin{table}[h]
        \centering
        \small
        \begin{tabular}{l c l}
            \toprule
            \textbf{Cor} & \textbf{Tupla (R, G, B)} & \textbf{Significado} \\
            \midrule
            Preto & \texttt{(0, 0, 0)} & Ausência total de luz \\
            Branco & \texttt{(255, 255, 255)} & Intensidade máxima de todos os canais \\
            Vermelho & \texttt{(255, 0, 0)} & Apenas canal R ativo \\
            Verde & \texttt{(0, 255, 0)} & Apenas canal G ativo \\
            Azul & \texttt{(0, 0, 255)} & Apenas canal B ativo \\
            Amarelo & \texttt{(255, 255, 0)} & Mistura de Vermelho + Verde \\
            Ciano & \texttt{(0, 255, 255)} & Mistura de Verde + Azul \\
            \bottomrule
        \end{tabular}
    \end{table}

    \begin{block}{Dica de Código}
        Defina cores como constantes no topo do arquivo com nomes em MAIÚSCULAS para deixar seu código legível: \texttt{AZUL = (50, 150, 250)}.
    \end{block}
\end{frame}

\begin{frame}[fragile]{Desenhando Primitivos com \texttt{pygame.draw}}
    O módulo \texttt{pygame.draw} permite criar formas na tela sem precisar de imagens externas:

    \begin{lstlisting}[language=Python]
# 1. Retangulo: (tela, cor, (x, y, largura, altura), espessura)
pygame.draw.rect(tela, (0, 120, 255), (100, 150, 80, 50))

# 2. Circulo: (tela, cor, (centro_x, centro_y), raio, espessura)
pygame.draw.circle(tela, (255, 50, 50), (400, 300), 40)

# 3. Linha: (tela, cor, (inicio_x, inicio_y), (fim_x, fim_y), espessura)
pygame.draw.line(tela, (255, 255, 255), (50, 50), (200, 50), 3)

# 4. Poligono (ex: triangulo): (tela, cor, [(x1,y1), (x2,y2), ...])
pygame.draw.polygon(tela, (255, 220, 0), [(300, 80), (250, 180), (350, 180)])
    \end{lstlisting}

    \vspace{0.2cm}
    \begin{itemize}
        \item Se o argumento de espessura (\texttt{width}) for omitido ou for $0$, a forma é \textbf{totalmente preenchida}.
        \item Se for maior que $0$, é desenhada apenas a \textbf{borda}.
    \end{itemize}
\end{frame}

\section{Animação, FPS e Interatividade}

\begin{frame}{Como Criar a Ilusão de Movimento?}
    \begin{columns}
        \begin{column}{0.5\textwidth}
            \textbf{A Lógica da Física:}
            \begin{itemize}
                \item Uma animação consiste em mudar levemente as coordenadas $(x, y)$ de um objeto a cada passagem pelo laço principal.
                \item Equação básica da velocidade constante:
                \[
                    x \leftarrow x + v_x
                \]
                \[
                    y \leftarrow y + v_y
                \]
            \end{itemize}
        \end{column}
        \begin{column}{0.5\textwidth}
            \begin{alertblock}{O Erro Clássico do Rastro}
                Se você esquecer de chamar \texttt{tela.fill(COR)} antes de desenhar o novo quadro, a tela não será limpa!
                \vspace{0.2cm}
                
                O objeto desenhado no quadro anterior continuará lá, criando um rastro borrado infinito.
            \end{alertblock}
        \end{column}
    \end{columns}
\end{frame}

\begin{frame}{O Papel Crucial de \texttt{pygame.time.Clock()}}
    \begin{block}{O que acontece sem o controle de FPS?}
        \begin{itemize}
            \item O computador rodará o \texttt{while} na velocidade máxima que o processador aguentar (milhares de ciclos por segundo).
            \item O jogo consumirá \textbf{100\% de CPU}, e rodará ultra rápido em máquinas potentes e lento em computadores antigos!
        \end{itemize}
    \end{block}

    \vspace{0.3cm}
    \begin{columns}
        \begin{column}{0.5\textwidth}
            \begin{block}{Com \texttt{relogio.tick(60)}}
                \begin{itemize}
                    \item O Pygame calcula quanto tempo o quadro demorou e coloca a CPU para dormir (\textit{sleep}) pelo tempo restante.
                    \item Garante que o jogo rode a exatos \textbf{60 FPS} em qualquer computador.
                \end{itemize}
            \end{block}
        \end{column}
        \begin{column}{0.5\textwidth}
            \textbf{Fórmula do Tempo de Quadro:}
            \[
                \Delta t = \frac{1000\text{ ms}}{60\text{ FPS}} \approx 16{,}6\text{ ms por quadro}
            \]
        \end{column}
    \end{columns}
\end{frame}

\begin{frame}[fragile]{Movimentação Interativa pelo Teclado}
    Para movimentar personagens de forma contínua e suave, utilizamos \texttt{pygame.key.get\_pressed()}:

    \begin{lstlisting}[language=Python]
# Captura o estado continuo de todas as teclas
teclas = pygame.key.get_pressed()

# Movimento Horizontal (Setas ou A / D)
if teclas[pygame.K_LEFT] or teclas[pygame.K_a]:
    jogador_x -= velocidade
if teclas[pygame.K_RIGHT] or teclas[pygame.K_d]:
    jogador_x += velocidade

# Movimento Vertical (Setas ou W / S)
if teclas[pygame.K_UP] or teclas[pygame.K_w]:
    jogador_y -= velocidade
if teclas[pygame.K_DOWN] or teclas[pygame.K_s]:
    jogador_y += velocidade
    \end{lstlisting}

    \vspace{0.2cm}
    \textbf{Diferença Fundamental:}
    \begin{itemize}
        \item \texttt{pygame.KEYDOWN} (na fila de eventos): Dispara \textbf{uma única vez} quando a tecla desce (ótimo para tiro ou pulo).
        \item \texttt{pygame.key.get\_pressed()}: Responde \textbf{enquanto a tecla estiver pressionada} (ótimo para corrida).
    \end{itemize}
\end{frame}

\section{Colisões com \texttt{pygame.Rect} e o Mini-Jogo}

\begin{frame}{O Poder da Classe \texttt{pygame.Rect}}
    O \textbf{\texttt{pygame.Rect}} é a estrutura mais versátil do Pygame para representar a posição e as dimensões de objetos 2D.

    \vspace{0.3cm}
    \begin{columns}
        \begin{column}{0.5\textwidth}
            \begin{block}{Criação}
                \texttt{r = pygame.Rect(x, y, larg, alt)}
            \end{block}
            \textbf{Propriedades automáticas:}
            \begin{itemize}
                \item \texttt{r.x}, \texttt{r.y}
                \item \texttt{r.width}, \texttt{r.height}
                \item \texttt{r.top}, \texttt{r.bottom}
                \item \texttt{r.left}, \texttt{r.right}
                \item \texttt{r.center}, \texttt{r.centerx}, \texttt{r.centery}
            \end{itemize}
        \end{column}
        \begin{column}{0.5\textwidth}
            \begin{exampleblock}{Bloqueio de Bordas Fácil}
                \begin{itemize}
                    \item \texttt{if r.left < 0: r.left = 0}
                    \item \texttt{if r.right > 800: r.right = 800}
                    \item \texttt{if r.top < 0: r.top = 0}
                    \item \texttt{if r.bottom > 600: r.bottom = 600}
                \end{itemize}
                \small{Não é necessário recalcular largura ou altura manualmente!}
            \end{exampleblock}
        \end{column}
    \end{columns}
\end{frame}

\begin{frame}[fragile]{Detecção de Colisão: \texttt{colliderect()}}
    Para saber se dois objetos se tocaram no espaço 2D, usamos a colisão por caixas alinhadas aos eixos (\textit{AABB - Axis-Aligned Bounding Box}):

    \vspace{0.2cm}
    \begin{lstlisting}[language=Python]
jogador = pygame.Rect(100, 100, 40, 40)
moeda   = pygame.Rect(300, 200, 20, 20)

# O metodo colliderect devolve True se houver sobreposicao!
if jogador.colliderect(moeda):
    pontos += 1
    # Reposiciona a moeda em outro local aleatorio
    moeda.x = random.randint(50, 750)
    moeda.y = random.randint(50, 550)
    \end{lstlisting}

    \vspace{0.2cm}
    \begin{block}{Princípio Matemático AABB}
        Dois retângulos colidem se, e somente se, houver sobreposição simultânea em ambos os eixos $X$ e $Y$:
        \[
            (A_{\text{left}} < B_{\text{right}}) \land (A_{\text{right}} > B_{\text{left}}) \land (A_{\text{top}} < B_{\text{bottom}}) \land (A_{\text{bottom}} > B_{\text{top}})
        \]
    \end{block}
\end{frame}

\begin{frame}[fragile]{Escrevendo Textos na Tela (HUD de Pontuação)}
    Para desenhar pontuações, vidas e mensagens na tela, seguimos três passos:

    \vspace{0.2cm}
    \begin{enumerate}
        \item \textbf{Carregar a Fonte:}
        \begin{lstlisting}[language=Python]
fonte = pygame.font.SysFont("Arial", 24, bold=True)
        \end{lstlisting}
        
        \item \textbf{Renderizar o Texto em uma Surface:}
        \begin{lstlisting}[language=Python]
# Argumentos: (texto, antialiasing_suave, cor_RGB)
img_texto = fonte.render(f"Pontos: {pontos}", True, (255, 255, 255))
        \end{lstlisting}

        \item \textbf{Transferir a Imagem para a Tela (\texttt{blit}):}
        \begin{lstlisting}[language=Python]
# Posiciona a imagem de texto nas coordenadas (x=20, y=20)
tela.blit(img_texto, (20, 20))
        \end{lstlisting}
    \end{enumerate}

    \begin{alertblock}{Atenção}
        A palavra \textbf{blit} vem de \textit{Block Image Transfer} (copiar blocos de pixels de uma superfície para outra).
    \end{alertblock}
\end{frame}

\section{Laboratório Prático e Conclusão}

\begin{frame}{Resumo dos Pilares Aprendidos}
    Nesta aula introdutória de 90 minutos, dominamos os quatro alicerces fundamentais da computação gráfica 2D:

    \vspace{0.3cm}
    \begin{columns}
        \begin{column}{0.5\textwidth}
            \begin{itemize}
                \item \textbf{Game Loop:} A espinha dorsal contínua que processa entradas, física e renderização.
                \item \textbf{Plano Cartesiano 2D:} Origem $(0, 0)$ no topo esquerdo e coordenada $Y$ invertida para baixo.
            \end{itemize}
        \end{column}
        \begin{column}{0.5\textwidth}
            \begin{itemize}
                \item \textbf{Controle de Tempo:} Uso de \texttt{Clock.tick(60)} para garantir taxa estável em qualquer hardware.
                \item \textbf{Colisões com Rect:} Otimização e praticidade através do método \texttt{colliderect()}.
            \end{itemize}
        \end{column}
    \end{columns}

    \vspace{0.4cm}
    \begin{block}{Boas Práticas de Engenharia de Jogos}
        \begin{itemize}
            \item Sempre limpe a tela com \texttt{tela.fill()} antes de desenhar.
            \item Trate o evento \texttt{pygame.QUIT} para não congelar o sistema operacional.
            \item Separe a lógica de atualização da renderização visual.
        \end{itemize}
    \end{block}
\end{frame}

\begin{frame}{Laboratório Prático: Desafios para a Turma}
    \begin{block}{Desafio 1: A Bola Multicolorida (Nível Fácil)}
        \begin{itemize}
            \item Execute o arquivo \texttt{desafio1\_bola\_quicando.py}.
            \item Faça a bola quicar continuamente nas 4 paredes da janela.
            \item \textbf{Bônus:} A cada batida, sorteie uma nova cor aleatória com \texttt{random.randint(50, 255)}.
        \end{itemize}
    \end{block}

    \vspace{0.2cm}
    \begin{block}{Desafio 2: Esquiva de Meteoros (Nível Intermediário)}
        \begin{itemize}
            \item Execute o arquivo \texttt{desafio2\_esquiva\_obstaculo.py}.
            \item Mova seu personagem na horizontal para desviar de blocos caindo do topo da tela.
            \item Utilize \texttt{colliderect()} para descontar vidas do jogador se for atingido!
        \end{itemize}
    \end{block}
\end{frame}

\begin{frame}{Próximos Passos no Curso}
    \begin{center}
        \Large{\textbf{O que exploraremos nas próximas aulas?}}
    \end{center}

    \vspace{0.3cm}
    \begin{itemize}
        \item \textbf{Sprites e Grupos:} Uso da classe \texttt{pygame.sprite.Sprite} e \texttt{pygame.sprite.Group} para gerenciar dezenas de entidades com POO.
        \item \textbf{Carregamento de Imagens:} Substituir formas geométricas por Spritesheets e animações de personagens.
        \item \textbf{Sistemas de Partículas e Câmera:} Efeitos visuais dinâmicos e mundos maiores que a tela.
    \end{itemize}

    \vspace{0.5cm}
    \begin{center}
        \textbf{Obrigado!} \\
        \footnotesize{IFCE Campus Tauá -- Curso de Análise e Desenvolvimento de Sistemas}
    \end{center}
\end{frame}


\end{document}
